problem:  
Let \((M^{2n+1}, T^{1,0}M, \theta)\) be a closed strictly pseudoconvex pseudohermitian manifold with Tanaka–Webster curvature tensor \(R_{\alpha\bar\beta}\) and torsion tensor \(A_{\alpha\beta}\). Denote by \(\Delta_b\) the sub-Laplacian and let \(\lambda_1(\Delta_b)\) be its first positive eigenvalue. Assume the pseudohermitian Ricci tensor satisfies, for some \(\kappa>0\),  
\[
\mathrm{Ric}_b(Z,\bar Z) \ge (n+1)\kappa\,|Z|_\theta^2 - \varepsilon\,|A|_\theta^2 |Z|_\theta^2
\quad \text{for all } Z\in T^{1,0}M,
\]
where \(|A|_\theta := \sup_{|Z|_\theta=1}|A_{\alpha\beta}Z^\alpha Z^\beta|\) and \(0 < \varepsilon \ll 1\).  

**Problem:**  
Is there a universal constant \(C_n>0\), depending only on the CR dimension \(n\), such that under the preceding curvature–torsion constraint one has the quantitative spectral bound  
\[
\lambda_1(\Delta_b) \ge n\kappa - C_n\,\varepsilon\,\|A\|_{L^\infty(M)}^2?
\]
Moreover, if a sequence of pseudohermitian structures \((\theta_\varepsilon)\) satisfies these conditions with \(\|A^{(\varepsilon)}\|_{L^\infty}\to 0\) and achieves asymptotic equality in the bound, does \((M, \theta_\varepsilon)\) converge (up to pseudohermitian scaling and CR automorphism) to the standard CR sphere \((S^{2n+1}, \theta_{\mathrm{std}})\)?  

Why is it a "good" problem:  
This problem asks whether the classical CR Lichnerowicz–Obata spectral gap is **quantitatively stable** under small pseudohermitian torsion perturbations. It uses a natural geometric control—an \(L^{\infty}\) smallness of the Tanaka–Webster torsion—and probes whether the fundamental eigenvalue of the sub-Laplacian retains a nearly extremal value \(n\kappa\) when the curvature bound is slightly weakened by torsion effects. The question combines **CR geometric analysis**, **spectral perturbation theory**, and **rigidity phenomena** from sub-Riemannian geometry. Despite its compact formulation, resolving it would require delicate analytical control of the Bochner–Weitzenböck identities in the pseudohermitian setting and new estimates showing how torsion enters the spectral geometry. Conceptually, it seeks a CR analogue of **stability of the Obata theorem** well known in Riemannian geometry. This simplicity, cross-disciplinary depth, and potential to clarify the stability of the CR sphere’s extremality make it a “good” and forward-looking research problem.